% GENERATED FILE. Edit canonical lessons, then run npm run generate:books.
% canonical-source-hash: fnv1a64:e1945887fada7692
% canonical-lessons: UR-C140-jitna, UR-C140-harna, UR-C140-sikhana, UR-C140-samjhana, UR-C140-mangna

\chapter{To win, To lose, To teach, To explain, To ask for}
\label{ch:ur-to-win-to-lose-to-teach-to-explain-to-ask-for}
\hlchaptermodality{140}

\begin{chapteropening}
\textbf{By the end of this chapter:} \emph{I can say to win, to lose, to teach, to explain and to ask for.}

Complete the last lesson of chapter 140: I can say to win, to lose, to teach, to explain and to ask for.
\end{chapteropening}

\section[\texorpdfstring{\ur{جیتنا} (jītnā)}{جیتنا (jītnā)}]{\ur{جیتنا} (jītnā) — to win}

\label{lesson:UR-C140-jitna}



\begin{quote}
Before the new one: say the Urdu for to run away, then the Urdu for to call.
\end{quote}

\subsection*{You'll want to know: \ur{جیتنا}}
\textbf{\ur{جیتنا}} — \emph{jītnā} — \textquotedblleft{}to win\textquotedblright{}.

\begin{cousinweb}
\textbf{\ur{جیتنا}} (\emph{jītnā}) is a tadbhava, a word inherited from Sanskrit and reshaped by sound change. It comes from \emph{jita}, \textquotedblleft{}won, conquered\textquotedblright{} (CDIAL 5224).
\end{cousinweb}

\begin{grammarlens}[title={What you've built}]
One more word for this chapter.
\end{grammarlens}

\subsection*{Your turn}
\begin{itemize}
\raggedright
  \item \emph{Say it:} \emph{jītnā}
  \item \emph{Say it:} \emph{jītnā}, once more
  \item \emph{Say it:} say \emph{jītnā}
  \item \emph{Recall:} say the Urdu for to run away, then the Urdu for to call, then say \emph{jītnā} again
\end{itemize}

\begin{tcolorbox}[breakable,colback=teal!4,colframe=teal!35!black,title={Before you move on}]
What does \ur{جیتنا} mean? (\textquotedblleft{}To win\textquotedblright{}.) Say it once more.
\end{tcolorbox}
\section[\texorpdfstring{\ur{ہارنا} (hārnā)}{ہارنا (hārnā)}]{\ur{ہارنا} (hārnā) — to lose (a game)}

\label{lesson:UR-C140-harna}



\begin{quote}
Before the new one: say the Urdu for to call, then the Urdu for to win.
\end{quote}

\subsection*{You'll want to know: \ur{ہارنا}}
\textbf{\ur{ہارنا}} — \emph{hārnā} — \textquotedblleft{}to lose (a game)\textquotedblright{}.

\begin{cousinweb}
\textbf{\ur{ہارنا}} (\emph{hārnā}) is a tadbhava, a word inherited from Sanskrit and reshaped by sound change. It comes from \emph{hārayati}, \textquotedblleft{}causes to be taken; loses\textquotedblright{} (CDIAL 14061).
\end{cousinweb}

\begin{grammarlens}[title={What you've built}]
One more word for this chapter.
\end{grammarlens}

\subsection*{Your turn}
\begin{itemize}
\raggedright
  \item \emph{Say it:} \emph{hārnā}
  \item \emph{Say it:} \emph{hārnā}, once more
  \item \emph{Say it:} say \emph{hārnā}
  \item \emph{Recall:} say the Urdu for to call, then the Urdu for to win, then say \emph{hārnā} again
\end{itemize}

\begin{tcolorbox}[breakable,colback=teal!4,colframe=teal!35!black,title={Before you move on}]
What does \ur{ہارنا} mean? (\textquotedblleft{}To lose (a game)\textquotedblright{}.) Say it once more.
\end{tcolorbox}
\section[\texorpdfstring{\ur{سکھانا} (sikhānā)}{سکھانا (sikhānā)}]{\ur{سکھانا} (sikhānā) — to teach}

\label{lesson:UR-C140-sikhana}



\begin{quote}
Before the new one: say the Urdu for to win, then the Urdu for to lose.
\end{quote}

\subsection*{You'll want to know: \ur{سکھانا}}
\textbf{\ur{سکھانا}} — \emph{sikhānā} — \textquotedblleft{}to teach\textquotedblright{}.

\begin{cousinweb}
\textbf{\ur{سکھانا}} (\emph{sikhānā}) is the causing form of \emph{sīkhnā}, \textquotedblleft{}to learn\textquotedblright{}: to teach is to make someone learn. CDIAL lists both under Sanskrit \emph{śikṣate}, \textquotedblleft{}learns\textquotedblright{} (CDIAL 12430).
\end{cousinweb}

\begin{grammarlens}[title={What you've built}]
One more word for this chapter.
\end{grammarlens}

\subsection*{Your turn}
\begin{itemize}
\raggedright
  \item \emph{Say it:} \emph{sikhānā}
  \item \emph{Say it:} \emph{sikhānā}, once more
  \item \emph{Say it:} say \emph{sikhānā}
  \item \emph{Recall:} say the Urdu for to win, then the Urdu for to lose, then say \emph{sikhānā} again
\end{itemize}

\begin{tcolorbox}[breakable,colback=teal!4,colframe=teal!35!black,title={Before you move on}]
What does \ur{سکھانا} mean? (\textquotedblleft{}To teach\textquotedblright{}.) Say it once more.
\end{tcolorbox}
\section[\texorpdfstring{\ur{سمجھانا} (samjhānā)}{سمجھانا (samjhānā)}]{\ur{سمجھانا} (samjhānā) — to explain}

\label{lesson:UR-C140-samjhana}



\begin{quote}
Before the new one: say the Urdu for to lose, then the Urdu for to teach.
\end{quote}

\subsection*{You'll want to know: \ur{سمجھانا}}
\textbf{\ur{سمجھانا}} — \emph{samjhānā} — \textquotedblleft{}to explain\textquotedblright{}.

\begin{cousinweb}
\textbf{\ur{سمجھانا}} (\emph{samjhānā}) is the causing form of \emph{samajhnā}, \textquotedblleft{}to understand\textquotedblright{}: to explain is to make someone understand. CDIAL lists the verb under Sanskrit \emph{saṃbudhyate}, \textquotedblleft{}wakes up, understands\textquotedblright{} (CDIAL 12959).
\end{cousinweb}

\begin{grammarlens}[title={What you've built}]
One more word for this chapter.
\end{grammarlens}

\subsection*{Your turn}
\begin{itemize}
\raggedright
  \item \emph{Say it:} \emph{samjhānā}
  \item \emph{Say it:} \emph{samjhānā}, once more
  \item \emph{Say it:} say \emph{samjhānā}
  \item \emph{Recall:} say the Urdu for to lose, then the Urdu for to teach, then say \emph{samjhānā} again
\end{itemize}

\begin{tcolorbox}[breakable,colback=teal!4,colframe=teal!35!black,title={Before you move on}]
What does \ur{سمجھانا} mean? (\textquotedblleft{}To explain\textquotedblright{}.) Say it once more.
\end{tcolorbox}
\section[\texorpdfstring{\ur{مانگنا} (māṅgnā)}{مانگنا (māṅgnā)}]{\ur{مانگنا} (māṅgnā) — to ask for}

\label{lesson:UR-C140-mangna}



\begin{quote}
Before the new one: say the Urdu for to teach, then the Urdu for to explain.
\end{quote}

\subsection*{You'll want to know: \ur{مانگنا}}
\textbf{\ur{مانگنا}} — \emph{māṅgnā} — \textquotedblleft{}to ask for\textquotedblright{}.

\begin{cousinweb}
\textbf{\ur{مانگنا}} (\emph{māṅgnā}) is a tadbhava, a word inherited from Sanskrit and reshaped by sound change. It comes from \emph{mārgati}, \textquotedblleft{}seeks, requests\textquotedblright{} (CDIAL 10074). It is built on \emph{mārga}, \textquotedblleft{}path\textquotedblright{}: to seek was \textquotedblleft{}to track\textquotedblright{}.
\end{cousinweb}

\begin{grammarlens}[title={What you've built}]
One more word for this chapter.
\end{grammarlens}

\subsection*{Your turn}
\begin{itemize}
\raggedright
  \item \emph{Say it:} \emph{māṅgnā}
  \item \emph{Say it:} \emph{māṅgnā}, once more
  \item \emph{Say it:} say \emph{māṅgnā}
  \item \emph{Recall:} say the Urdu for to teach, then the Urdu for to explain, then say \emph{māṅgnā} again
\end{itemize}

\begin{tcolorbox}[breakable,colback=teal!4,colframe=teal!35!black,title={Before you move on}]
What does \ur{مانگنا} mean? (\textquotedblleft{}To ask for\textquotedblright{}.) Say it once more.
\end{tcolorbox}
